\documentclass[11pt]{article}

\usepackage[margin=0.9in]{geometry}
\usepackage{amsmath,amssymb}
\usepackage{booktabs}
\usepackage{enumitem}
\usepackage{listings}
\usepackage{xcolor}
\usepackage{hyperref}
\usepackage{tikz}

\definecolor{backcolour}{rgb}{0.95,0.95,0.92}
\definecolor{codeblue}{rgb}{0.05,0.15,0.55}
\definecolor{codegreen}{rgb}{0.0,0.45,0.0}

\lstset{
    backgroundcolor=\color{backcolour},
    basicstyle=\ttfamily\small,
    keywordstyle=\color{codeblue}\bfseries,
    commentstyle=\color{codegreen},
    stringstyle=\color{purple},
    breaklines=true,
    frame=single,
    showstringspaces=false,
    tabsize=4
}

\title{CPSC 250 \\ Craps Game Odds}
\author{Edward Brash}
\date{}

\begin{document}

\maketitle

\section{Motivation}

The craps-style dice game used in the class example has an interesting property:

\begin{quote}
The probability of winning one credit is exactly equal to the probability of losing one credit.
\end{quote}

That is:

\[
P(+1) = P(-1) = \frac{1}{2}
\]

This means the game behaves like a random walk.

\section{Game Rules}

On the initial roll:

\begin{itemize}
    \item rolling 7 or 11 wins one credit;
    \item rolling 3 or 12 loses one credit;
    \item any other roll sets a goal and requires more rolling.
\end{itemize}

On later rolls:

\begin{itemize}
    \item rolling 7 loses one credit;
    \item rolling the goal wins one credit;
    \item anything else means roll again.
\end{itemize}

\section{Initial Roll Probabilities}

The number of possible outcomes when rolling two dice is:

\[
36
\]

The probabilities of common sums are:

\[
P(7) = \frac{6}{36} = \frac{1}{6}
\]

\[
P(11) = \frac{2}{36} = \frac{1}{18}
\]

\[
P(3) = \frac{2}{36} = \frac{1}{18}
\]

\[
P(12) = \frac{1}{36}
\]

Therefore:

\[
P(7 \text{ or } 11) = \frac{1}{6} + \frac{1}{18}
= \frac{4}{18}
= \frac{2}{9}
\]

and:

\[
P(3 \text{ or } 12) = \frac{1}{18} + \frac{1}{36}
= \frac{3}{36}
= \frac{1}{12}
\]

\section{Re-Roll Cases}

If the initial roll is neither an immediate win nor an immediate loss, the game enters a re-roll state.

For the simplified version in the notes, this leads to a repeated process:

\begin{itemize}
    \item some rolls end the round,
    \item some rolls continue the round.
\end{itemize}

The repeated continuation creates a geometric series.

\section{Geometric Series Review}

A geometric series has the form:

\[
a + ar + ar^2 + ar^3 + \cdots
\]

If \(|r| < 1\), then:

\[
a + ar + ar^2 + ar^3 + \cdots = \frac{a}{1-r}
\]

Equivalently:

\[
\sum_{n=0}^{\infty} ar^n = \frac{a}{1-r}
\]

\section{Probability of Eventually Winning}

The notes derive the probability of winning as a sum involving immediate wins plus delayed wins.

The structure is:

\[
P(+1) = \text{immediate win} + \text{delayed win}
\]

The delayed portion includes the possibility of rolling again and again before eventually ending the round.

This repeated structure is exactly why a geometric series appears.

\section{Final Result}

After simplifying the geometric series, the result is:

\[
P(+1) = \frac{1}{2}
\]

Similarly:

\[
P(-1) = \frac{1}{2}
\]

Since every completed round either wins or loses one credit,

\[
P(+1) + P(-1) = 1
\]

and therefore, once one probability is known to be \(1/2\), the other must also be \(1/2\).

\section{Interpretation}

The game is fair in the sense that:

\[
P(+1) = P(-1)
\]

But this does not mean the game lasts forever.

The player's credits move up and down like a random walk until eventually reaching zero.

\section{Random Walk Model}

If the player begins with some number of credits, then after each round:

\[
\text{credits} \rightarrow \text{credits} + 1
\]

with probability \(1/2\), or:

\[
\text{credits} \rightarrow \text{credits} - 1
\]

with probability \(1/2\).

The game ends when:

\[
\text{credits} = 0
\]

\section{Simulation Philosophy}

The notes then transition from probability to simulation.

Instead of only solving analytically, we can simulate the game many times.

The idea is:

\begin{enumerate}
    \item Run the game to completion many times.
    \item Record the number of rounds in each game.
    \item Record credits as a function of round number.
    \item Use those results for analysis and plotting.
\end{enumerate}

\section{Data Structures for Simulation}

For many simulations, we might use lists such as:

\begin{lstlisting}[language=Python]
credit_result = []
round_result = []
result_number = []
\end{lstlisting}

For each simulated game, we append:

\begin{itemize}
    \item the list of credit values,
    \item the list of round numbers,
    \item the final number of rounds.
\end{itemize}

\section{Simulation Loop}

A simplified structure is:

\begin{lstlisting}[language=Python]
num_sims = 1000

credit_result = []
round_result = []
result_number = []

for i in range(num_sims):

    credit_round = []
    result_round = []

    # run one complete game here

    round_result.append(result_round)
    credit_result.append(credit_round)
    result_number.append(rounds)
\end{lstlisting}

\section{Why This Matters}

This example connects several important ideas:

\begin{itemize}
    \item probability,
    \item loops,
    \item lists of lists,
    \item simulations,
    \item plotting,
    \item random walks.
\end{itemize}

It also gives a real reason to store data during a simulation.

\section{Key Takeaways}

\begin{itemize}
    \item The dice game has equal probabilities of winning or losing one credit.
    \item Repeated re-rolls naturally lead to geometric series.
    \item The game behaves like a random walk.
    \item Simulations can be used to study behavior that is difficult to understand from a single run.
    \item Lists are useful for storing simulation histories.
\end{itemize}

\end{document}
